\section{Tamamlayıcı vaka: hangi ortalama, kimin oranı?}
\label{sec:v4-star98}

\textbf{Soru ve kapsam.} STAR98'in aynı sayımlarıyla iki soru soralım:
``Dosyadaki bölgelerin oran ortalaması nedir?'' ve ``Bu sayımlarda temsil
edilen öğrencilerin toplam üst kategori oranı nedir?'' Yeni bir regresyon
kurmayacağız; özetlenen hedefin neden değiştiğini inceleyeceğiz.

Dağıtıcının açıklaması satırları birleşik okul bölgesi, yanıtı 9. sınıf
matematik sonuçlarının sayımı olarak tanımlar; dosya özgün kaynağın alt
kümesidir. Aynı sayfanın notlarında geçen \textit{counties} ifadesiyle
adlandırma çelişkisi çözülmüş değildir. Burada ``bölge'', açıklamadaki
tanıma dayanır; gerçek coğrafi kimlikler bağımsız doğrulanmış sayılmaz.\footnote{Dağıtıcı tanımı:
\url{https://www.statsmodels.org/stable/datasets/generated/star98.html};
erişim 13 Eylül 2026. Kaynak eser: \citet{gill2001}.}

\textbf{Analiz tanımı.} Yerel 303 satırın tamamı kullanılır. $a_j$, kaynakta
\texttt{NABOVE}; $b_j$, \texttt{NBELOW}; $n_j=a_j+b_j$ ve $r_j=a_j/n_j$
olsun. Üst kategori kaynağın ulusal medyan üstü tanımıdır; medyana eşit
puanların işlemine ilişkin yeni bir kural varsayılmaz. Payda bütün kayıtlı
öğrenciler değil, bu iki sayımın toplamıdır. Kimlik yerel sıra anahtarıdır.
Eksik, negatif, kesirli sayım veya sıfır payda varsa işlem durur; otomatik
silme ya da doldurma yapılmaz. Sıfır üst sayımı ise geçerlidir.

\subsection*{İki hedefi hesaplamadan önce tanımlayın}

$J=303$ için eşit bölge ağırlıklı özet
\[
\bar r_{\mathrm{bölge}}=\frac{1}{J}\sum_{j=1}^{J}r_j
\]
her satıra $1/J$ ağırlığı verir. Toplam sayım oranı ise
\[
r_{\mathrm{toplam}}=\frac{\sum_j a_j}{\sum_j n_j}
=\sum_j\frac{n_j}{\sum_k n_k}r_j
\]
olup satırları sayım büyüklükleriyle ağırlıklandırır. Her iki formül de
aynı kayıtların betimlemesidir; ikinci formülün ağırlıkları ters dahil
edilme olasılığı değildir. Büyük bölgeye daha fazla ağırlık vermek,
kendiliğinden örnekleme yanlılığını düzeltmez.

\begin{table}[H]
\centering
\small
\caption{V4: aynı 303 kayıtta iki farklı hedef}
\label{tab:v4-hedefler}
\begin{tabularx}{\linewidth}{@{}Xr@{}}
\toprule
Özet & Değer \\
\midrule
Üst kategori toplamı & 108\,418 \\
İki kategorinin toplamı & 267\,611 \\
Eşit bölge ağırlıklı oran (yüzde) & 43,697840 \\
Toplam sayım oranı (yüzde) & 40,513282 \\
Toplam eksi bölge özeti (yüzde puan) & $-3{,}184558$ \\
\bottomrule
\end{tabularx}
\end{table}

Kayıt başına toplam 33 ile 38\,852 arasında değişir. Bu veride büyüklükle
ağırlıklandırılmış oran daha düşüktür. Bu fark bir yazılım uyuşmazlığı
değildir; hedef değişmiştir. Her büyük bölgenin oranının her küçük
bölgeden düşük olduğu sonucu da çıkmaz. B11--B12'de soru değişince
ortalama ve oran ayrılıyordu; burada değişken aynı kalırken ağırlıklandırma
hedefi değişmektedir.

\noindent\begin{minipage}{\linewidth}
\textbf{Yeniden üretim.} Proje kökünden, yalnız betimsel iki özeti hesaplayın:
\begin{lstlisting}[style=prismPythonApp]
import pandas as pd

data = pd.read_csv(
    "companion/data/raw/star98_statsmodels.csv")
counts = data[["NABOVE", "NBELOW"]]
assert len(data) == 303
assert data["DISTRICT_ID"].is_unique
assert counts.notna().all().all()
assert (counts >= 0).all().all()
assert (counts % 1 == 0).all().all()
total = counts.sum(axis=1)
assert (total > 0).all()
rate = counts["NABOVE"] / total
equal_rate = rate.mean()
pooled_rate = counts["NABOVE"].sum() / total.sum()
print(100 * equal_rate, 100 * pooled_rate)
print(100 * (pooled_rate - equal_rate))
\end{lstlisting}
\end{minipage}

Kaynak hash'i, yerel anahtar ve ham--temiz aktarım denetimi dahil tam kod
\path{companion/vakalar/v4-star98/analyze.py} dosyasındadır.
Bu özetler için güven aralığı veya \pval-değeri üretilmez: dosyanın
betimlenmesi, özgün örnekleme mekanizmasının ya da öğrencilerin bağımsız
Bernoulli gözlemleri olduğunun doğrulanması değildir. 267\,611 sayısını
bağımsız örneklem büyüklüğü sayıp otomatik binom aralığı hesaplamayın.

\subsection*{Toplulaştırma hangi bilgiyi kaybettirir?}

Aşağıdaki iki tablo \textbf{tamamen yapaydır}; STAR98'in bireyleri değildir.
Her senaryoda 100 kişi, 50 düşük gelirli kişi ve toplam 50 üst kategori
sonucu vardır. Yine de grup içi ilişkiler ters yöndedir.

\begin{table}[H]
\centering
\small
\caption{Aynı marjinleri veren iki yapay bireysel düzen}
\label{tab:v4-yapay}
\begin{tabular}{@{}llrr@{}}
\toprule
Senaryo & Grup & Üst & Alt \\
\midrule
A & Düşük gelir & 40 & 10 \\
A & Diğer & 10 & 40 \\
B & Düşük gelir & 10 & 40 \\
B & Diğer & 40 & 10 \\
\bottomrule
\end{tabular}
\end{table}

İki senaryoda da düşük gelir oranı ve üst kategori oranı \%50'dir. Fakat
düşük gelir eksi diğer grup oran farkı A'da $0{,}8-0{,}2=0{,}6$,
B'de $0{,}2-0{,}8=-0{,}6$ olur. Yalnız marjinleri bilmek bireysel
ilişkiyi belirlemez. Bölümdeki bölgesel regresyon katsayısını bireysel gelir
etkisi sayamayız. Bu gösterim gerçek STAR98'de bir terslenme kanıtı değildir;
toplulaştırmanın hangi bilgiyi sakladığını gösterir. B13'ün öğrenci düzeyi
uygulamasında birim farklıdır; veri dosyaları birleştirilmez.

\subsection{Görevler, kontrol yanıtları ve puanlama}

Her görev iki bileşenlidir; her bileşen 1 puan, toplam \textbf{10 puan}.
Önce yanıtınızı yazın, sonra aşağıdaki gerekçelerle karşılaştırın.
\begin{enumerate}
\item Bir satırın birimini ve toplam sayım oranının paydasını tanımlayın.
  267\,611 sayısı hangi evreni otomatik olarak temsil etmez?
\item İki oranın neden farklı olduğunu ağırlıklarıyla açıklayın.
  ``Büyük olan doğru, küçük olan hatalı'' yorumunu değerlendirin.
\item $n_j/\sum_k n_k$ ağırlığının örnekleme ağırlığı olup olmadığını
  tartışın. Toplam sayımdan neden hemen bir binom güven aralığı üretmiyoruz?
\item Yapay A ve B'de düşük gelir eksi diğer grup oran farkını bulun.
  Aynı marjinlerden bireysel ilişkinin neden belirlenemediğini açıklayın.
\item İki gerçek özeti içeren, genelleme ve nedensellik sınırları açık
  kısa bir rapor yazın. Yapay kişileri gerçek toplamın dışında tutun.
\end{enumerate}

\textbf{Gerekçeli kontrol yanıtları.}
\begin{enumerate}
\item Satır, dağıtıcı açıklamasındaki bölge düzeyi toplulaştırmadır;
  payda iki kayıtlı sonuç sayısının toplamıdır (1). Kaynak alt kümesi,
  bütün California/ABD öğrencilerinin evreni değildir; idari adlandırma
  çelişkisi ayrıca korunur (1).
\item Birinci özet her satıra $1/303$, ikinci özet sayım büyüklüğüyle
  orantılı ağırlık verir (1). Hedefler farklı olduğundan fark hata veya
  üstünlük kanıtı değildir; rapor soruya uygun hedefi belirtmelidir (1).
\item Bunlar büyüklük ağırlıklarıdır; dahil edilme olasılıkları
  belgelenmiş değildir (1). Bireysel bağımsızlık ve örnekleme modeli
  kurulmadığından toplam sayı doğrudan binom belirsizliğini haklı kılmaz (1).
\item A'da $+0{,}6$, B'de $-0{,}6$ (60 ve eksi 60 yüzde puan) bulunur (1).
  Aynı marjinlerle zıt ilişkiler mümkün olduğundan bireysel ilişki
  belirlenemez; bunun gerçek terslenme kanıtı olmadığı belirtilir (1).
\item Rapor iki özeti doğru adlarıyla verir ve yapay tabloları gerçek
  hesaplara katmaz (1). Örnekleme kapsamını aşan genelleme ve bireysel
  nedensellik kurmaz (1); doğru sayılar tek başına tam puan getirmez.
\end{enumerate}

\textbf{Kısa rapor örneği.} ``Yerel STAR98 dosyasındaki 303 toplulaştırılmış
kayıtta eşit bölge ağırlıklı üst kategori oranı ortalaması \%43,697840;
108\,418 üst sayımın 267\,611 toplam sayıma oranı \%40,513282'dir.
Fark ağırlıklandırma hedefinden kaynaklanır. Sonuçlar bu kayıtları betimler;
tüm öğrenciler için evren tahmini veya bireysel gelir etkisi değildir.
Örnekleme belirsizliği bu betimsel karşılaştırmada hesaplanmamıştır.''